\clearpage
\setcounter{chapter}{3}
\chaptermark{介质界面光学与近场光学显微镜}
\pdfbookmark[0]{第三章 介质界面光学与近场光学显微镜}{chapter3}
\problem{3-1}{
光矢量与入射面之夹角被简称为振动的方位角或偏振角。设入射的线偏振光的方位角为 $\alpha_1$，而入射角为 $i$。试证明，反射线偏振光的方位角 $\alpha'_1$ 和折射线偏振光的方位角 $\alpha_2$ 分别由以下两式给出，
\[
\tan\alpha'_1=-\frac{\cos(i_2-i_1)}{\cos(i_2+i_1)}\tan\alpha_1,\qquad
\tan\alpha_2=\frac{n_2\cos i_1+n_1\cos i_2}{n_1\cos i_1+n_2\cos i_2}\tan\alpha'_1.
\]
}
\problem{3-2}{
一束线偏振光从空气入射到玻璃表面上，其入射角恰巧为布儒斯特角，而方位角为 $20^\circ$。试求反射线偏振和折射线偏振的方位角 $\alpha'_1$ 和 $\alpha_2$，设玻璃折射率为 $1.56$。
}
\problem{3-3}{
试计算：
\begin{enumerate}
\item 光从空气入射于水面的布儒斯特角 $i_{\mathrm B}$，水的折射率为 $4/3$。
\item 一束自然光从水入射于某种玻璃表面上，当入射角为 $50.82^\circ$ 时反射光成为线偏振光，该玻璃的折射率为多少？
\end{enumerate}
}
\problem{3-4}{
设入射光、反射光和折射光的总光功率分别为 $W_1$、$W'_1$ 和 $W_2$，则总光功率的反射率 $\mathcal R$ 和透射率 $\mathcal T$ 定义为
\[
\mathcal R=\frac{W'_1}{W_1},\qquad \mathcal T=\frac{W_2}{W_1}.
\]
下面的问题均系 $(\mathcal R,\mathcal T)$ 与 $(\mathcal R_{\mathrm p},\mathcal T_{\mathrm p})$、$(\mathcal R_{\mathrm s},\mathcal T_{\mathrm s})$ 的关系。
\begin{enumerate}
\item 当入射光为线偏振光且其方位角为 $\alpha$ 时，试证明
\[
\mathcal R=\mathcal R_{\mathrm p}\cos^2\alpha+\mathcal R_{\mathrm s}\sin^2\alpha,\qquad
\mathcal T=\mathcal T_{\mathrm p}\cos^2\alpha+\mathcal T_{\mathrm s}\sin^2\alpha.
\]
\item 当入射光为自然光时，试证明
\[
\mathcal R=\frac12(\mathcal R_{\mathrm p}+\mathcal R_{\mathrm s}),\qquad
\mathcal T=\frac12(\mathcal T_{\mathrm p}+\mathcal T_{\mathrm s}).
\]
\item 当入射光为圆偏振光时，试证明
\[
\mathcal R=\frac12(\mathcal R_{\mathrm p}+\mathcal R_{\mathrm s}),\qquad
\mathcal T=\frac12(\mathcal T_{\mathrm p}+\mathcal T_{\mathrm s}).
\]
\end{enumerate}
}
\problem{3-5}{
一线偏振光以 $45^\circ$ 角入射于一玻璃面，其方位角为 $60^\circ$，玻璃折射率为 $1.50$。求
\begin{enumerate}
\item 光功率反射率 $\mathcal R$ 和透射率 $\mathcal T$。
\item 若改为自然光入射，$\mathcal R$ 和 $\mathcal T$ 变为多少？
\end{enumerate}
}
\problem{3-6}{
如图所示，一束自然光入射于一平板玻璃，现观测到反射光强 $I_1=0.1I_0$。求
\begin{enumerate}
\item 图中标出的各光束 $2$、$3$、$4$，其光功率 $W_2$、$W_3$、$W_4$ 为多少？设最初入射光功率为 $W_0$，并忽略吸收。
\item 若要求出光强比 $I_2/I_0$，还应当给出什么条件？
\end{enumerate}
\tikzfig{fig_3_6}{3-6}
}
\problem{3-7}{
在光于介质表面的反射和折射实验中，获得以下测量数据：入射角 $i_1\approx75^\circ$，折射角 $i_2\approx40^\circ$，总光强反射率 $R\approx30\%$，试求出，
\begin{enumerate}
\item 总振幅反射率 $r$，总光功率反射率 $\mathcal R$；
\item 总光功率透射率 $\mathcal T$，总光强透射率 $T$ 和总振幅透射率 $t$。
\item 这入射光是自然光吗？它可能是何种偏振态？
\end{enumerate}
}
\problem{3-8}{
如图所示，它是一个由半导体材料砷化镓（GaAs）制成的发光管，其管芯 $(AB)$ 为发光区，直径 $d$ 约 $3\,\mathrm{mm}$。为了避免全反射，发光管上部被研磨成半球形，以使管芯发出的光有最大的透射率向外发射。若要求发光区周边 $A$、$B$ 发的光不发生全反射，那半球的半径 $r$ 应当为多少？已知 GaAs 的折射率为 $3.4$，对其发光的光波波长为 $0.9\,\mu\mathrm m$。
\tikzfig{fig_3_8}{3-8}
}
\problem{3-9}{
计算玻片组的偏振度。取用三块平板玻璃叠放一起，其折射率为 $1.64$。当一束自然光以布儒斯特角 $i_{\mathrm B}$ 入射时，
\begin{enumerate}
\item 最终从这玻片组透射出来的光其偏振度 $p$ 为多少？
\item 其中所含 $\mathrm s$ 光的强度 $I_{\mathrm s}$ 与强度 $I_{\mathrm p}$ 之比值为多少？
\end{enumerate}
}
\problem{3-10}{
考察过临界角时全反射光的相位变化情况。设一光束由玻璃入射于空气，且入射角已超过临界角 $i_{\mathrm c}$，玻璃折射率为 $1.68$。
\begin{enumerate}
\item 取用一张坐标纸来绘制反射 $\mathrm s$ 光和 $\mathrm p$ 光的相移 $\delta'_{\mathrm s}(i)$ 曲线和 $\delta'_{\mathrm p}(i)$ 曲线，要求横坐标即入射角 $i$ 在 $i_{\mathrm c}$—$\pi/2$ 范围中至少选取 $6$ 个数据点。
\item 利用这两条曲线找出相移差值 $\delta(i)=\delta'_{\mathrm p}(i)-\delta'_{\mathrm s}(i)$ 的最大值 $\delta_{\mathrm M}$ 和相应的入射角 $i_{\mathrm M}$。
\item 设入射光为一线偏振光，其方位角在一、三象限即 $\varphi_{1\mathrm s}=\varphi_{1\mathrm p}=0$，当入射角为 $60^\circ$ 时，问反射光的两个正交振动之间的相位差 $\delta=\varphi'_{1\mathrm s}-\varphi'_{1\mathrm p}$ 为多少？其合成的椭圆偏振光是左旋的还是右旋的？
\item 当入射角为 $45^\circ$ 或 $75^\circ$ 时，在透射区即空气中存在的隐失波的穿透深度 $d_1$ 或 $d_2$ 为多少？设光波长为 $633\,\mathrm{nm}$。
\end{enumerate}
}
